\section{Trace a match without confusing its bits with its output}
The fixed-match fixture produces 35 zero bytes. Its complete zlib stream is ten bytes:
two wrapper bytes, four compressed bytes and four Adler bytes. The compressed fields
occupy exactly 32 bits in this fixture.

\Needspace{12\baselineskip}
\begin{longtable}{p{40mm}p{22mm}p{26mm}p{29mm}p{27mm}}
\toprule Boundary & New bits & Total bits & Next input offset & Output bytes \\\midrule
Final fixed header & 3 & 3 & 3 & 0 \\
Literal zero & 8 & 11 & 4 & 1 \\
Length symbol 272 & 7 & 18 & 5 & 1 \\
Length extra field 3 & 2 & 20 & 5 & 1 \\
Distance symbol zero & 5 & 25 & 6 & 1 \\
Zero distance extras & 0 & 25 & 6 & 1 \\
Forward match copied & 0 & 25 & 6 & 35 \\
End-of-block & 7 & 32 & 6 & 35 \\\bottomrule
\end{longtable}
Input offsets include the two zlib header bytes. A field can consume bits without
advancing the next-byte cursor, and the history copy produces 34 bytes without consuming
compressed bits. After end-of-block, the retained compressed cursor is at trailer
offset six and the output cursor is at 35. The exact checks and checksum comparison
can now follow the same path as a stored block.

\section{Compare three descriptions of the same dynamic lengths}
The three repeat fixtures reconstruct identical alphabets. Literal zero has length
one, end-of-block and length symbol 257 have length two, and distance symbol zero has
length one. Zero lengths between those present entries describe absent symbols.
The following checkpoints occur before the prepared repeat has been stored.

\Needspace{9\baselineskip}
\begin{longtable}{p{27mm}p{25mm}p{27mm}p{24mm}p{36mm}}
\toprule Operation & Cursor before & Extra field & Zero run & Cursor after store \\\midrule
16 & 2 & 3 in two bits & 6 & 8 \\
17 & 1 & 7 in three bits & 10 & 11 \\
18 & 1 & 127 in seven bits & 138 & 139 \\\bottomrule
\end{longtable}
The rest of each description fills the combined range to 259 entries. At completion,
indices 0, 256, 257 and 258 contain 1, 2, 2 and 1; indices 1 through 255 contain zero.
Different encoded descriptions can therefore produce exactly the same tables and
filtered output. This is one reason to grade behavior rather than one source listing.

Try the malformed and mixed recipes after all regions are complete. A distance before
available history, reserved data symbols, an initial previous-length repeat, an invalid
prefix shape, a missing end code or a bad Adler must fail with zero returned count.
The mixed stored/fixed input must succeed. Inspect scratch rows separately from final
context outputs and framebuffer words: a failed call may leave a produced prefix in
scratch, but the temporary whole program never publishes pixels.

For a later optimization exercise, consider precomputing immutable fixed tables instead
of reconstructing them for each fixed block. Keep dynamic tables private to the current
call and retain every stream, history and output check. Compare identical output and
measured instruction work including setup. The measurements elsewhere in this chapter
retain their original sources and boundaries; they are not new results for this
construction sequence.
